\documentclass[11pt]{article}
\usepackage[margin=1in]{geometry}
\usepackage{enumitem}
% =============================================================================
% Preambles/header.tex — shared LaTeX/Beamer preamble for this project
%
% Use from a Beamer lecture by prepending:
%     \input{header}
% and compiling with TEXINPUTS=../Preambles:$TEXINPUTS (the /compile-latex
% skill does this automatically).
%
% PALETTE CONTRACT. Color names defined here MUST match the names in
% Quarto/theme-template.scss so Beamer and Quarto renderings use the same
% palette. The scripts/check-palette-sync.sh script enforces this.
%
% To customize: edit the HEX values below AND the matching $variable in
% Quarto/theme-template.scss, then run scripts/check-palette-sync.sh.
% =============================================================================

% -----------------------------------------------------------------------------
% Engine detection + encoding
%
% Our /compile-latex skill uses XeLaTeX, where inputenc is unnecessary and
% can produce encoding warnings. Only load inputenc under pdfTeX.
% -----------------------------------------------------------------------------
\usepackage{iftex}
\ifPDFTeX
  \usepackage[utf8]{inputenc}
\fi

\usepackage{amsmath, amssymb, amsthm}
\usepackage{booktabs}
\usepackage{graphicx}

% -----------------------------------------------------------------------------
% PALETTE — mirrors Quarto/theme-template.scss variable names.
% Load xcolor and define colors BEFORE anything that references them
% (notably hyperref, hypersetup, and the Beamer color assignments).
% -----------------------------------------------------------------------------
\usepackage{xcolor}

\definecolor{primary-blue}{HTML}{012169}      % headings, accents
\definecolor{primary-gold}{HTML}{B9975B}      % emphasis, borders
\definecolor{highlight-yellow}{HTML}{F2A900}  % markers, alerts
\definecolor{light-bg}{HTML}{E8EDF5}          % title / section backgrounds
\definecolor{jet}{HTML}{1A1A1A}               % body text

% Semantic colors (used in diagrams and annotations)
\definecolor{positive}{HTML}{15803D}          % good / observed / identified
\definecolor{negative}{HTML}{B91C1C}          % bad / problematic / bias
\definecolor{neutral}{HTML}{525252}           % reference / context

% Highlight accents (match .hi-* classes in the SCSS)
\definecolor{hi-slate}{HTML}{314F4F}
\definecolor{hi-green}{HTML}{2E8B57}
\definecolor{hi-red}{HTML}{C41E3A}

% Semantic aliases so skill templates and snippets can write
% \textcolor{accent}{...} without hard-coding "primary-blue".
\colorlet{accent}{primary-blue}
\colorlet{accent2}{primary-gold}

% -----------------------------------------------------------------------------
% Hyperref — loaded AFTER xcolor + palette so link colors resolve.
% -----------------------------------------------------------------------------
\usepackage{hyperref}
\hypersetup{colorlinks=true, linkcolor=primary-blue, urlcolor=primary-blue,
            citecolor=primary-blue, pdfborder={0 0 0}}

% -----------------------------------------------------------------------------
% Course code — set once per deck (after \input{header}, before \begin{document})
% via \coursecode{CS401}. Shown in the footer on every slide so a deck's course
% is identifiable without opening the title page. Empty by default.
% -----------------------------------------------------------------------------
\makeatletter
\newcommand{\coursecode}[1]{\gdef\@coursecodetext{#1}}
\gdef\@coursecodetext{}
\makeatother

% -----------------------------------------------------------------------------
% Beamer theme assignments (only apply under Beamer).
%
% We detect Beamer via \@ifclassloaded{beamer}, which is the documented,
% non-internal way. This must be wrapped in \makeatletter/\makeatother
% because \@ifclassloaded contains an @-catcode character.
% -----------------------------------------------------------------------------
\makeatletter
\@ifclassloaded{beamer}{%
  \usefonttheme{professionalfonts}
  \setbeamercolor{structure}{fg=primary-blue}
  \setbeamercolor{frametitle}{fg=primary-blue}
  \setbeamercolor{title}{fg=primary-blue}
  \setbeamercolor{subtitle}{fg=primary-gold}
  \setbeamercolor{author}{fg=primary-blue}
  \setbeamercolor{institute}{fg=neutral}
  \setbeamercolor{date}{fg=primary-gold}
  \setbeamercolor{itemize item}{fg=primary-blue}
  \setbeamercolor{itemize subitem}{fg=primary-gold}
  \setbeamercolor{alerted text}{fg=primary-gold}
  \setbeamercolor{block title}{fg=white, bg=primary-blue}
  \setbeamercolor{block body}{bg=light-bg}
  % Semantic block variants: block=definition/concept (blue), exampleblock=
  % worked example (green/positive), alertblock=key takeaway or caution (gold).
  % Keep to <=2 colored boxes per slide across ALL three variants (INV-7).
  \setbeamercolor{block title example}{fg=white, bg=positive}
  \setbeamercolor{block body example}{bg=light-bg}
  \setbeamercolor{block title alerted}{fg=white, bg=primary-gold}
  \setbeamercolor{block body alerted}{bg=light-bg}

  % Minimal footer: course code (left) + page X / N (right), no navigation clutter
  \setbeamertemplate{navigation symbols}{}
  \setbeamertemplate{footline}{%
    \color{neutral}\footnotesize\hspace{0.7em}\@coursecodetext\hfill
    \insertframenumber/\inserttotalframenumber\hspace{0.7em}\vspace{0.3em}}
}{}
\makeatother

% -----------------------------------------------------------------------------
% TikZ setup — libraries the snippet gallery uses
% -----------------------------------------------------------------------------
\usepackage{tikz}
\usetikzlibrary{arrows.meta, positioning, calc, decorations.pathreplacing,
                fit, shapes.geometric, backgrounds}

% Shared TikZ styles — snippets and hand-written diagrams can pick these up
% via `\begin{tikzpicture}[every node/.style={...}]` or by naming specific
% styles (e.g., `\node[dag-node] (X) at (0,0) {X};`).
\tikzset{
  % Node styles for causal DAGs and similar
  dag-node/.style = {
    draw=primary-blue, fill=light-bg, rounded corners=2pt,
    minimum width=1.2cm, minimum height=0.9cm, align=center, font=\small
  },
  decision-node/.style = {
    draw=primary-gold, fill=white, diamond, aspect=1.8,
    minimum width=1.8cm, minimum height=1.2cm, align=center, font=\small,
    inner sep=1pt
  },
  % Edge styles — observed vs counterfactual
  observed-edge/.style = {-{Latex[length=2.5mm]}, thick, draw=jet},
  counterfactual-edge/.style = {-{Latex[length=2.5mm]}, thick, draw=neutral, dashed},
  confound-edge/.style = {-{Latex[length=2.5mm]}, thick, draw=negative, dashed},
  % Data-point styles for plots
  observed-dot/.style = {fill=primary-blue, draw=primary-blue, circle, inner sep=1.2pt},
  counterfactual-dot/.style = {fill=white, draw=neutral, circle, inner sep=1.2pt}
}

% -----------------------------------------------------------------------------
% Convenience macros
% -----------------------------------------------------------------------------
\newcommand{\muted}[1]{\textcolor{neutral}{#1}}
\newcommand{\key}[1]{\textcolor{primary-gold}{\textbf{#1}}}
\newcommand{\good}[1]{\textcolor{positive}{#1}}
\newcommand{\bad}[1]{\textcolor{negative}{#1}}

% Standout frame macro: full-bleed dark background for section transitions.
% Beamer-only (uses \begin{frame}/\setbeamercolor/\usebeamerfont) -- do not
% call from an article-class file (e.g. Notes/<CODE>/*-notes.tex). \muted,
% \key, \good, \bad, and \coursecode above are plain xcolor/gdef and are
% safe to use from either class; verified when Notes/article-class support
% was added.
% Expand: \transitionslide{Your transition title}
\newcommand{\transitionslide}[1]{%
  \begingroup
  \setbeamercolor{background canvas}{bg=primary-blue}
  \begin{frame}[plain]
    \centering
    \vfill
    {\usebeamerfont{title}\color{white}\Large #1\par}
    \vfill
  \end{frame}
  \endgroup
}

% Section divider frame (Beamer-only), an alternative to \transitionslide with a
% darker two-line style: near-black (jet) background, a small white label line
% above a larger gold title, and a thin gold rule along the bottom edge.
% Expand: \sectiondivider{Section 2}{Pointers}
\newcommand{\sectiondivider}[2]{%
  \begingroup
  \setbeamercolor{background canvas}{bg=jet}
  \begin{frame}[plain]
    \vspace*{3.0cm}
    {\color{white}\ttfamily\large #1\par}
    \vspace{0.5cm}
    {\color{primary-gold}\LARGE #2\par}
    \begin{tikzpicture}[remember picture, overlay]
      \draw[primary-gold, line width=2.5pt]
        ($(current page.south west)+(0,0.1)$) -- ($(current page.south east)+(0,0.1)$);
    \end{tikzpicture}
  \end{frame}
  \endgroup
}

% -----------------------------------------------------------------------------
% Channel watermark — "Concepts That Click" (YouTube).
%
% Article-class files (CompetitiveExam Q/A sets, Notes, Assignments) get a
% light diagonal draft watermark on every page plus a centered footer naming
% the channel. Beamer decks get a subtle TikZ background watermark instead
% (the footline already carries the course code). Centralized here so every
% MCQ PDF is branded without per-file edits.
% -----------------------------------------------------------------------------
\makeatletter
\@ifclassloaded{article}{%
  \usepackage{draftwatermark}
  \SetWatermarkText{Concepts That Click}
  \SetWatermarkScale{1}
  \SetWatermarkFontSize{2cm}
  \SetWatermarkLightness{0.86}
  \SetWatermarkAngle{45}
  \usepackage{fancyhdr}
  \pagestyle{fancy}
  \fancyhf{}
  \fancyfoot[C]{\footnotesize\textcolor{neutral}{Concepts That Click~|~YouTube: \href{https://www.youtube.com/@concepts.thatclick}{@concepts.thatclick}}}
  \renewcommand{\headrulewidth}{0pt}
  \renewcommand{\footrulewidth}{0pt}
  \fancypagestyle{plain}{%
    \fancyhf{}%
    \fancyfoot[C]{\footnotesize\textcolor{neutral}{Concepts That Click~|~YouTube: \href{https://www.youtube.com/@concepts.thatclick}{@concepts.thatclick}}}%
    \renewcommand{\headrulewidth}{0pt}%
    \renewcommand{\footrulewidth}{0pt}%
  }
}{}
\@ifclassloaded{beamer}{%
  \setbeamertemplate{background}{%
    \begin{tikzpicture}[remember picture, overlay]
      \node[rotate=45, opacity=0.055, align=center]
        at (current page.center)
        {\Huge\textcolor{neutral}{Concepts That Click}};
    \end{tikzpicture}%
  }
}{}
\makeatother 
\coursecode{JKSSB}
\title{Lesson 59 Practice: Current CPU Architecture, Data Representation \& RISC/CISC}
\author{JKSSB Computer Awareness}
\date{\today}
\begin{document}
\maketitle
\noindent\textbf{Provenance:} 17 original JKSSB-pattern questions and five
paraphrased adaptations of the JKSSB Junior Assistant 2026 computer block
(Q62, Q64, Q65, final key). The adaptations at Q1, Q5, Q7, Q11, and Q13 are
reworded, not verbatim reproductions.

\begin{enumerate}[label=\textbf{Q\arabic*.}, leftmargin=*]
\item \textit{[PYQ-adapted: JKSSB Junior Assistant 2026 Q62; final key;
paraphrased]} Which statement describes von Neumann (stored-program)
architecture?
  \begin{enumerate}[label=\Alph*.]
    \item Instructions and data are stored in two separate memories.
    \item Only data is stored in memory; instructions are not.
    \item Instructions and data are stored together in the same memory.
    \item Only instructions are stored in memory; data is not.
  \end{enumerate}
\item \textit{[Original]} In immediate addressing, the operand is:
  \begin{enumerate}[label=\Alph*.]
    \item held in a general-purpose register.
    \item contained within the instruction itself.
    \item stored in main memory at an address named by the instruction.
    \item the address of the next instruction to fetch.
  \end{enumerate}
\item \textit{[Original]} The cache-hit ratio is defined as:
  \begin{enumerate}[label=\Alph*.]
    \item hits divided by total accesses.
    \item misses divided by total accesses.
    \item hits divided by misses.
    \item total accesses divided by hits.
  \end{enumerate}
\item \textit{[Original]} Why is a cache hit faster than a cache miss?
  \begin{enumerate}[label=\Alph*.]
    \item The cache is located on secondary storage.
    \item The cache stores only instructions, never data.
    \item The cache is smaller and faster and sits closer to the CPU, so its
      access time is lower than main memory.
    \item The cache runs at the same speed as a hard disk.
  \end{enumerate}
\item \textit{[PYQ-adapted: JKSSB Junior Assistant 2026 Q62; final key;
paraphrased]} Which statement about data hazards is correct?
  \begin{enumerate}[label=\Alph*.]
    \item A hazard can arise only when an operand is in memory.
    \item Register addressing removes every possible hazard.
    \item A hazard can arise only with immediate addressing.
    \item A hazard can still arise from a dependency between instructions,
      for example in a pipeline.
  \end{enumerate}
\item \textit{[Original]} Which option correctly contrasts cache and virtual
memory?
  \begin{enumerate}[label=\Alph*.]
    \item Cache is on secondary storage, while virtual memory is in the
      cache.
    \item Cache and virtual memory are the same piece of hardware.
    \item Virtual memory is a single cache line.
    \item Cache is small fast memory between the CPU and main memory;
      virtual memory makes memory appear larger by using secondary storage as
      backing store.
  \end{enumerate}
\item \textit{[PYQ-adapted: JKSSB Junior Assistant 2026 Q64; final key;
paraphrased]} The Program Counter holds:
  \begin{enumerate}[label=\Alph*.]
    \item the address of the instruction currently being executed.
    \item the address of the next instruction to be fetched.
    \item the operand value of the current instruction.
    \item the last instruction that finished executing.
  \end{enumerate}
\item \textit{[Original]} The Instruction Register holds:
  \begin{enumerate}[label=\Alph*.]
    \item the address of the next instruction to fetch.
    \item the instruction currently being fetched, decoded, and executed.
    \item the value of the operand.
    \item the width of the data bus.
  \end{enumerate}
\item \textit{[Original]} Interrupts are normally processed:
  \begin{enumerate}[label=\Alph*.]
    \item only after the entire program has completed.
    \item only during system boot.
    \item only for hardware faults.
    \item at instruction boundaries, during program execution.
  \end{enumerate}
\item \textit{[Original]} Which statement about instruction timing is correct?
  \begin{enumerate}[label=\Alph*.]
    \item Every machine instruction takes exactly one clock cycle.
    \item The clock-cycle count is fixed for the whole program.
    \item Different machine instructions can take different numbers of clock
      cycles.
    \item Machine instructions do not consume clock cycles.
  \end{enumerate}
\item \textit{[PYQ-adapted: JKSSB Junior Assistant 2026 Q65; final key;
paraphrased]} In two's-complement arithmetic, overflow occurs when:
  \begin{enumerate}[label=\Alph*.]
    \item there is a carry out of the most significant bit.
    \item the result of an addition is exactly zero.
    \item the operation is performed on unsigned numbers.
    \item two operands of the same sign produce a result of the opposite
      sign.
  \end{enumerate}
\item \textit{[Original]} When widening a value to more bits, which rule is
correct?
  \begin{enumerate}[label=\Alph*.]
    \item Zero-extension is for signed values and sign-extension is for
      unsigned values.
    \item Zero-extension is for unsigned values and sign-extension is for
      signed values.
    \item Both unsigned and signed values use zero-extension.
    \item Both unsigned and signed values use sign-extension.
  \end{enumerate}
\item \textit{[PYQ-adapted: JKSSB Junior Assistant 2026 Q65; final key;
paraphrased]} In BCD (binary-coded decimal), the decimal number 25 is
represented as:
  \begin{enumerate}[label=\Alph*.]
    \item the plain binary value $11001_2$.
    \item eight bits per decimal digit.
    \item the hexadecimal value $25_{16}$.
    \item $0010\ 0101$, one 4-bit code for each decimal digit.
  \end{enumerate}
\item \textit{[Original]} Which statement about byte-addressability and the
effective address is correct?
  \begin{enumerate}[label=\Alph*.]
    \item The instruction length is always equal to the effective address.
    \item Memory is word-addressable only, so individual bytes have no
      address.
    \item Each byte has its own address, and the effective address is the
      operand's actual location, which need not equal the instruction
      length.
    \item The effective address is always equal to the Program Counter value.
  \end{enumerate}
\item \textit{[Original]} What are the full forms of RISC and CISC?
  \begin{enumerate}[label=\Alph*.]
    \item RISC = Reduced Instruction Set Computer; CISC = Complex
      Instruction Set Computer
    \item RISC = Rapid Instruction Set Computer; CISC = Complex Instruction
      Set Computer
    \item RISC = Reduced Integer Storage Cache; CISC = Central Instruction
      Set Computer
    \item RISC = Random Instruction Set Computing; CISC = Common Instruction
      Set Computer
  \end{enumerate}
\item \textit{[Original]} Which set of features describes a RISC processor?
  \begin{enumerate}[label=\Alph*.]
    \item Many complex, variable-length instructions.
    \item Instructions that always operate directly on memory operands.
    \item Microprogrammed control of highly complex instructions.
    \item Simple, mostly fixed-length instructions, many registers, and a
      load/store architecture.
  \end{enumerate}
\item \textit{[Original]} Which set of features describes a CISC processor?
  \begin{enumerate}[label=\Alph*.]
    \item A few simple fixed-length instructions only.
    \item A large set of complex, variable-length instructions, often
      allowing memory operands.
    \item No general-purpose registers at all.
    \item A strict load/store-only memory access rule.
  \end{enumerate}
\item \textit{[Original]} Which statement is \textbf{NOT} correct?
  \begin{enumerate}[label=\Alph*.]
    \item The Program Counter holds the address of the next instruction to
      fetch.
    \item Interrupts can be handled at instruction boundaries.
    \item All machine instructions take the same number of clock cycles.
    \item Two's-complement overflow can occur when two same-sign operands are
      added.
  \end{enumerate}
\item \textit{[Original]} Which code correctly represents the decimal digit 5
in BCD?
  \begin{enumerate}[label=\Alph*.]
    \item $0101$
    \item $1001$
    \item $101$
    \item $1110$
  \end{enumerate}
\item \textit{[Original]} Evaluate the statements:
  \begin{enumerate}[label=\Roman*.]
    \item In immediate addressing, the operand is part of the instruction.
    \item The cache-hit ratio equals hits divided by total accesses.
    \item Two's-complement overflow is the same as a carry out of the most
      significant bit.
  \end{enumerate}
  \begin{enumerate}[label=\Alph*.]
    \item I and II only
    \item II and III only
    \item I and III only
    \item I, II, and III
  \end{enumerate}
\item \textit{[Original]} The von Neumann bottleneck refers to:
  \begin{enumerate}[label=\Alph*.]
    \item the delay between keyboard input and screen output.
    \item the limited rate at which instructions and data can be transferred
      between the CPU and memory over a shared path.
    \item the time taken for a hard disk to spin up.
    \item the number of general-purpose registers in the CPU.
  \end{enumerate}
\item \textit{[Original]} The main purpose of virtual memory is to:
  \begin{enumerate}[label=\Alph*.]
    \item increase the number of CPU registers.
    \item replace the cache entirely.
    \item reduce the CPU clock speed.
    \item make the available memory appear larger than physical RAM by using
      secondary storage as backing store.
  \end{enumerate}
\end{enumerate}
\end{document}
